% ============================================================================
% CHAPTER 1C: ELECTRICAL SYSTEMS
% ============================================================================

\chapter{Electrical Systems}
\label{ch:electrical}

Electrical circuits provide clean examples of dynamic systems with well-defined mathematical models. The concepts developed here---impedance, frequency response, filtering---form the historical foundation of control theory.
\begin{historicalbox}[title=Historical Context: Black's Negative Feedback Amplifier and Frequency Response Methods]
The foundation of modern control theory emerged not from mechanical engineers but from electrical engineers at Bell Telephone Laboratories. In 1927, Harold Black invented the negative feedback amplifier, recognizing that feeding a portion of the output back inverted to the input could stabilize amplifier gain against component variations and temperature changes. This elegant insight---that feedback could improve performance---was initially rejected by patent offices, who considered it obvious. Yet Black's feedback amplifier enabled long-distance telephone transmission by compensating for signal loss over cables, fundamentally changing telecommunications.

The mathematical analysis of feedback systems was subsequently developed by Harry Nyquist (1932) and Hendrik Bode (1945), who created graphical methods to analyze stability and performance from frequency response measurements. The Nyquist stability criterion and Bode plots transformed control design from an empirical art into a quantitative science. These frequency response techniques, born from electrical circuits and amplifiers, became universal tools applicable to all domains---mechanical, thermal, biological, and beyond---establishing electrical systems as the historical birthplace of modern control theory.
\end{historicalbox}

% ============================================================================
\section{Passive Circuit Elements}
\label{sec:passive_elements}
% ============================================================================

\subsection{Resistor}

Ohm's law relates voltage and current:
\begin{equation}
v(t) = Ri(t)
\end{equation}

In the Laplace domain: $V(s) = RI(s)$

Impedance: $Z_R(s) = R$

\input{figures/fig_ch09_resistor.tex}

\subsection{Capacitor}

The capacitor stores charge:
\begin{equation}
i(t) = C\frac{\dd v}{\dd t} \quad \Leftrightarrow \quad v(t) = \frac{1}{C}\int i(\tau)\,\dd\tau
\end{equation}

In the Laplace domain: $I(s) = CsV(s)$ or $V(s) = \frac{1}{Cs}I(s)$

Impedance: $Z_C(s) = \frac{1}{Cs}$

\input{figures/fig_ch09_capacitor.tex}

\subsection{Inductor}

The inductor stores magnetic energy:
\begin{equation}
v(t) = L\frac{\dd i}{\dd t} \quad \Leftrightarrow \quad i(t) = \frac{1}{L}\int v(\tau)\,\dd\tau
\end{equation}

In the Laplace domain: $V(s) = LsI(s)$ or $I(s) = \frac{1}{Ls}V(s)$

Impedance: $Z_L(s) = Ls$

\input{figures/fig_ch09_inductor.tex}

\input{figures/fig_ch09_impedance_plane.tex}

% ============================================================================
\section{RC Circuit (First-Order System)}
\label{sec:rc_circuit}
% ============================================================================

\input{figures/fig_ch09_rc_circuit.tex}

\subsection{Time-Domain Analysis}

Applying KVL:
\begin{equation}
v_{\text{in}} = v_R + v_{\text{out}} = Ri + v_{\text{out}}
\end{equation}

Since $i = C\frac{\dd v_{\text{out}}}{\dd t}$:
\begin{equation}
\boxed{RC\frac{\dd v_{\text{out}}}{\dd t} + v_{\text{out}} = v_{\text{in}}}
\end{equation}

This is a first-order ODE with time constant $\tau = RC$.

\subsection{Transfer Function}

\begin{equation}
G(s) = \frac{V_{\text{out}}(s)}{V_{\text{in}}(s)} = \frac{1}{RCs + 1} = \frac{1}{\tau s + 1}
\end{equation}

\subsection{Frequency Response}

Substituting $s = \jj\omega$:
\begin{equation}
G(\jj\omega) = \frac{1}{1 + \jj\omega RC}
\end{equation}

Magnitude and phase:
\begin{align}
|G(\jj\omega)| &= \frac{1}{\sqrt{1 + (\omega RC)^2}} \\
\angle G(\jj\omega) &= -\arctan(\omega RC)
\end{align}

Cutoff frequency (where magnitude is $1/\sqrt{2}$ or $-3$ dB):
\begin{equation}
\omega_c = \frac{1}{RC} = \frac{1}{\tau}
\end{equation}

\begin{examplebox}[title=Example 1.6: RC Circuit Step Response]
An RC circuit has $R = 10$ k$\Omega$ and $C = 1$ $\mu$F. Find:
\begin{enumerate}[label=(\alph*)]
    \item Time constant
    \item Cutoff frequency
    \item Response to a 5V step input
\end{enumerate}

\textbf{Solution:}

(a) Time constant:
\begin{equation}
\tau = RC = 10{,}000 \times 10^{-6} = 0.01 \text{ s} = 10 \text{ ms}
\end{equation}

(b) Cutoff frequency:
\begin{equation}
f_c = \frac{1}{2\pi RC} = \frac{1}{2\pi \times 0.01} = 15.9 \text{ Hz}
\end{equation}

(c) Step response:
\begin{equation}
v_{\text{out}}(t) = 5(1 - \ee^{-t/0.01}) = 5(1 - \ee^{-100t}) \text{ V}
\end{equation}

At $t = \tau = 10$ ms: $v_{\text{out}} = 5(1 - \ee^{-1}) = 3.16$ V (63.2\% of final)
\end{examplebox}

\input{figures/fig_ch09_rc_step.tex}

\input{figures/fig_ch09_rc_bode.tex}

% ============================================================================
\section{RL Circuit}
\label{sec:rl_circuit}
% ============================================================================

\input{figures/fig_ch09_rl_circuit.tex}

For current as output:
\begin{equation}
L\frac{\dd i}{\dd t} + Ri = v_{\text{in}}
\end{equation}

Transfer function (voltage to current):
\begin{equation}
\frac{I(s)}{V_{\text{in}}(s)} = \frac{1}{Ls + R} = \frac{1/R}{(L/R)s + 1}
\end{equation}

Time constant: $\tau = L/R$

% ============================================================================
\section{RLC Circuit (Second-Order System)}
\label{sec:rlc_circuit}
% ============================================================================

\input{figures/fig_ch09_rlc_circuit.tex}

\subsection{Differential Equation}

Applying KVL:
\begin{equation}
v_{\text{in}} = v_R + v_L + v_C = Ri + L\frac{\dd i}{\dd t} + v_{\text{out}}
\end{equation}

Since $i = C\frac{\dd v_{\text{out}}}{\dd t}$:
\begin{equation}
\boxed{LC\frac{\dd^2 v_{\text{out}}}{\dd t^2} + RC\frac{\dd v_{\text{out}}}{\dd t} + v_{\text{out}} = v_{\text{in}}}
\end{equation}

\subsection{Transfer Function}

\begin{equation}
G(s) = \frac{V_{\text{out}}(s)}{V_{\text{in}}(s)} = \frac{1}{LCs^2 + RCs + 1}
\end{equation}

In standard form:
\begin{equation}
G(s) = \frac{\omega_n^2}{s^2 + 2\zeta\omega_n s + \omega_n^2}
\end{equation}

where:
\begin{align}
\omega_n &= \frac{1}{\sqrt{LC}} \quad \text{(resonant frequency)} \\
\zeta &= \frac{R}{2}\sqrt{\frac{C}{L}} \quad \text{(damping ratio)}
\end{align}

\subsection{Quality Factor}

The quality factor $Q$ measures the sharpness of resonance:
\begin{equation}
Q = \frac{1}{2\zeta} = \frac{1}{R}\sqrt{\frac{L}{C}}
\end{equation}

High $Q$ means sharp resonance peak and low damping.

\begin{examplebox}[title=Example 1.7: RLC Circuit Design]
Design an RLC circuit with resonant frequency $f_n = 1$ kHz and $Q = 10$.

\textbf{Solution:}

Choose $C = 0.1$ $\mu$F. Then:
\begin{equation}
L = \frac{1}{(2\pi f_n)^2 C} = \frac{1}{(2\pi \times 1000)^2 \times 10^{-7}} = 0.253 \text{ H}
\end{equation}

For $Q = 10$:
\begin{equation}
R = \frac{1}{Q}\sqrt{\frac{L}{C}} = \frac{1}{10}\sqrt{\frac{0.253}{10^{-7}}} = 159 \text{ }\Omega
\end{equation}

Damping ratio: $\zeta = 1/(2Q) = 0.05$ (highly underdamped)
\end{examplebox}

\input{figures/fig_ch09_rlc_frequency_response.tex}

% ============================================================================
\section{Operational Amplifier Circuits}
\label{sec:opamp}
% ============================================================================

Operational amplifiers (op-amps) are fundamental building blocks for analog control systems.

\subsection{Ideal Op-Amp Assumptions}

\begin{itemize}
    \item Infinite open-loop gain: $A \to \infty$
    \item Infinite input impedance: $Z_{\text{in}} \to \infty$ (no input current)
    \item Zero output impedance: $Z_{\text{out}} = 0$
    \item Infinite bandwidth
\end{itemize}

With negative feedback, these lead to the ``golden rules'':
\begin{enumerate}
    \item $V_+ = V_-$ (virtual short)
    \item $I_+ = I_- = 0$ (no input current)
\end{enumerate}

\begin{figure}[htbp]
\centering
\begin{circuitikz}[scale=1.2]
    \draw (0,0) node[op amp] (opamp) {}
    (opamp.+) node[left] {$V_+$}
    (opamp.-) node[left] {$V_-$}
    (opamp.out) node[right] {$V_{\text{out}}$};
\end{circuitikz}
\caption{Op-amp symbol.}
\label{fig:opamp_symbol}
\end{figure}

\subsection{Inverting Amplifier}

\begin{figure}[htbp]
\centering
\begin{circuitikz}[scale=1]
    \draw (0,0) node[op amp] (opamp) {};
    \draw (opamp.+) -- ++(-0.5,0) node[ground] {};
    \draw (opamp.-) -- ++(-0.5,0) coordinate (A)
          to[R, l=$R_1$] ++(-2,0) node[left] {$V_{\text{in}}$};
    \draw (A) -- ++(0,1.5) coordinate (B)
          to[R, l=$R_2$] ++(2.5,0) coordinate (C)
          -- (C |- opamp.out) -- (opamp.out);
    \draw (opamp.out) -- ++(0.5,0) node[right] {$V_{\text{out}}$};
\end{circuitikz}
\caption{Inverting amplifier.}
\label{fig:inverting_amp}
\end{figure}

Transfer function:
\begin{equation}
\frac{V_{\text{out}}}{V_{\text{in}}} = -\frac{R_2}{R_1}
\end{equation}

\subsection{Non-Inverting Amplifier}

Transfer function:
\begin{equation}
\frac{V_{\text{out}}}{V_{\text{in}}} = 1 + \frac{R_2}{R_1}
\end{equation}

\subsection{Integrator}

\begin{figure}[htbp]
\centering
\begin{circuitikz}[scale=1]
    \draw (0,0) node[op amp] (opamp) {};
    \draw (opamp.+) -- ++(-0.5,0) node[ground] {};
    \draw (opamp.-) -- ++(-0.5,0) coordinate (A)
          to[R, l=$R$] ++(-2,0) node[left] {$V_{\text{in}}$};
    \draw (A) -- ++(0,1.5) coordinate (B)
          to[C, l=$C$] ++(2.5,0) coordinate (C)
          -- (C |- opamp.out) -- (opamp.out);
    \draw (opamp.out) -- ++(0.5,0) node[right] {$V_{\text{out}}$};
\end{circuitikz}
\caption{Op-amp integrator.}
\label{fig:integrator}
\end{figure}

Transfer function:
\begin{equation}
G(s) = \frac{V_{\text{out}}(s)}{V_{\text{in}}(s)} = -\frac{1}{RCs}
\end{equation}

Time-domain:
\begin{equation}
v_{\text{out}}(t) = -\frac{1}{RC}\int_0^t v_{\text{in}}(\tau)\,\dd\tau
\end{equation}

\subsection{Differentiator}

Transfer function:
\begin{equation}
G(s) = \frac{V_{\text{out}}(s)}{V_{\text{in}}(s)} = -RCs
\end{equation}

\begin{warningbox}
Pure differentiators amplify high-frequency noise. In practice, add a small resistor in series with $C$ to limit high-frequency gain.
\end{warningbox}

\subsection{PID Controller Implementation}

A PID controller can be implemented with op-amps:

\begin{equation}
G_{PID}(s) = K_p + \frac{K_i}{s} + K_d s
\end{equation}

\input{figures/fig_ch09_pid_block.tex}

% ============================================================================
\section{Active Filters}
\label{sec:active_filters}
% ============================================================================

\subsection{First-Order Low-Pass Filter}

\begin{equation}
G(s) = \frac{K}{s/\omega_c + 1}
\end{equation}

Passes low frequencies, attenuates high frequencies.

\subsection{First-Order High-Pass Filter}

\begin{equation}
G(s) = \frac{Ks/\omega_c}{s/\omega_c + 1}
\end{equation}

Passes high frequencies, attenuates low frequencies (including DC).

\subsection{Second-Order Butterworth Low-Pass}

Maximally flat passband response:
\begin{equation}
G(s) = \frac{\omega_c^2}{s^2 + \sqrt{2}\omega_c s + \omega_c^2}
\end{equation}

Damping ratio $\zeta = 1/\sqrt{2} \approx 0.707$, giving no overshoot in step response.

\subsection{Band-Pass Filter}

\begin{equation}
G(s) = \frac{K(\omega_0/Q)s}{s^2 + (\omega_0/Q)s + \omega_0^2}
\end{equation}

Passes frequencies near $\omega_0$, attenuates both lower and higher frequencies.

\input{figures/fig_ch09_filter_comparison.tex}

% ============================================================================
\section{Power Electronics}
\label{sec:power_electronics}
% ============================================================================

\subsection{Buck Converter (Step-Down)}

\input{figures/fig_ch09_buck.tex}

Steady-state output (ideal):
\begin{equation}
V_{\text{out}} = D \cdot V_{\text{in}}
\end{equation}

where $D$ is the duty cycle ($0 \leq D \leq 1$).

Small-signal transfer function (control to output):
\begin{equation}
\frac{\hat{v}_{\text{out}}(s)}{\hat{d}(s)} = V_{\text{in}} \cdot \frac{1}{LCs^2 + (L/R_L)s + 1}
\end{equation}

% ============================================================================
